% Generated draft: requires mathematical and source review.
\documentclass[11pt]{article}
\usepackage[margin=25mm]{geometry}
\usepackage{fontspec}
\setmainfont{Latin Modern Roman}
\usepackage{amsmath,amssymb,mathtools,unicode-math}
\setmathfont{Latin Modern Math}
\usepackage{enumitem,needspace,xurl,hyperref}
\setlength{\parindent}{0pt}
\setlength{\parskip}{5pt}
\setlength{\emergencystretch}{3em}
\newcommand{\N}{\mathbb{N}}
\newcommand{\Z}{\mathbb{Z}}
\newcommand{\Q}{\mathbb{Q}}
\newcommand{\R}{\mathbb{R}}
\newcommand{\C}{\mathbb{C}}
\newcommand{\F}{\mathbb{F}}
\newcommand{\A}{\mathbb{A}}
\newcommand{\PP}{\mathbb{P}}
\newcommand{\E}{\mathbb{E}}
\newcommand{\eps}{\varepsilon}
\newcommand{\dd}{\,\mathrm{d}}
\begin{document}
\section*{UnsolvedMath ID: 2307017}
\textbf{Source problem code:} catalogue entry
\subsection{Definitions and mathematical statement}
Let $H^+=\{z\in\C:\operatorname{Re}z>0\}$, let $f$ be analytic and bounded on $H^+$, and fix $\alpha\in\R$ with $|\alpha|<\pi/2$. The whole ray $re^{i\alpha}$, $r>0$, lies in $H^+$. With the signed-logarithm directional-type convention, define
\[
 h_f(\alpha)=\limsup_{r\to\infty}\frac{\log|f(re^{i\alpha})|}{r},
\]
using $\log0=-\infty$. For a sample sequence tending to infinity define
\[
 h_f((r_n),\alpha)=\limsup_{n\to\infty}
                   \frac{\log|f(r_ne^{i\alpha})|}{r_n}.
\]
These may be negative or $-\infty$ despite boundedness; replacing $\log$ by $\log^+$ would lose the intended decay information.

The source assumes positive integer sample radii with
\[
 \sum_{n=1}^{\infty}\frac1{r_n}=\infty
\]
and asks to prove
\[
 h_f((r_n),\alpha)=h_f(\alpha).
\]
It does not explicitly require that $r_n$ increase, be distinct, or tend to infinity. Those conditions are essential to interpreting sampling type and the density hypothesis: repetitions can make the reciprocal sum diverge without providing additional sampling radii. A standard conditional intended version takes $r_1<r_2<\cdots$ positive integers and retains the divergent reciprocal sum, but that is not silently substituted for the literal statement. The source also does not define its exponential-type convention, so the meaningful signed-logarithm reading above should be confirmed.

The source attributes proofs for $\alpha=0$ to Boas and Levinson and says their methods do not seem to work when $|\alpha|>\pi/4$. This is a reported limitation of those proofs, not a statement that all intermediate angles are settled.
\subsection{Short English statement}
Does a sufficiently dense integer-radius sample recover the full exponential decay/growth type of a bounded analytic function along any ray inside the right half-plane? The precise intended sampling condition needs confirmation because the printed sequence assumptions allow repetitions. The supplied triage reports the question open in Hayman's 2018 update and no definitive resolution located; no independent status review is asserted.
\subsection{Sources}
\begin{itemize}
\item UnsolvedMath permanent ID 2307017, source code catalogue entry, ``Research Problems in Function Theory --- Problem 7.17.'' Registry curation and metadata: Ulam AI and UnsolvedMath Contributors, CC BY 4.0; updated 2026-08-06.
\item Hayman, \emph{Research Problems in Function Theory} (2018 source update), Problem 7.17, attributed to J. Korevaar: \url{https://arxiv.org/abs/1809.07200}. Supplied source-TeX extraction accessed 2026-07-29. The metadata retains the update for status triage and marks status/release and rights review NEEDS\_REVIEW.
\item Boas and Levinson are named for proofs at angle zero without complete titles, dates or locators. No further theorem or reference is inferred.
\end{itemize}
\end{document}
